\documentclass[11pt,a4paper]{article}
\usepackage[margin=2.2cm]{geometry}
\usepackage{fontspec}
\setmainfont{TeX Gyre Pagella}
\usepackage{enumitem}
\usepackage{xcolor,listings}
\lstset{basicstyle=\ttfamily\small,frame=single,rulecolor=\color{gray!50},columns=fullflexible,keepspaces=true,breaklines=true}
\usepackage[hidelinks]{hyperref}
\setlist{nosep}
\setlength{\parskip}{4pt}

\title{Competitive Programming and Contests\\Hands-On 1 -- Report}
\author{Gabriele Marsili}
\date{October 2026}

\begin{document}\sloppy
\maketitle

The solution is in \texttt{lib.rs}: I added two public methods to the provided \texttt{Tree},
\texttt{is\_bst} and \texttt{sum\_max\_simple\_path}, and a set of tests.
In the following, $n$ is the number of nodes and $h$ the height of the tree.

\section*{Exercise 1 -- Checking whether the tree is a BST}

Considering the definition of a BST, it is necessary to decide a convention for keys with the same value,
otherwise it is not clear how duplicates are handled: I decided to allow duplicates in the right subtree of
a node.

An interesting observation is that comparing a node only with its children is not enough, because the
subtrees rooted at the children could contain keys in the wrong order. For this reason I introduced an
interval of admissible values for the key, passed down in the recursion (line 100):
\begin{lstlisting}
fn rec_is_bst(&self, node_id: Option<usize>, admissible_interval: (u32, u32)) -> bool
\end{lstlisting}
This handles cases like the test \texttt{grandchild\_violating\_ancestor\_is\_rejected}: with root 10, left
child 5 and 15 as the right child of 5, every parent-child pair is ordered, but 15 is in the left subtree
of 10.

The visit is recursive. If the id of a node is \texttt{None} the node does not exist and I return
\texttt{true} (base case). Otherwise I first check that the key is in the admissible interval and then I make
the recursive calls on the subtrees (inductive cases) with the intervals $[\mathit{low}, \mathit{key}-1]$ for
the left one and $[\mathit{key}, \mathit{high}]$ for the right one. At the root the interval is
$[\texttt{u32::MIN}, \texttt{u32::MAX}]$.

\textbf{Observation:} keys are \texttt{u32}, so $\mathit{key}-1$ would underflow when the key is 0. I use
\texttt{checked\_sub} (line 114), which returns \texttt{None} in this case: since no \texttt{u32} key is
smaller than 0, the left subtree must be empty.

\textbf{Complexity:}
\begin{itemize}
\item $O(n)$ time, since each node is visited at most once;
\item $O(h)$ extra space for the recursion stack, where $h$ is the height of the tree: $h = \log n$ if the
tree is balanced, $h = n$ in the worst case (a chain).
\end{itemize}

\section*{Exercise 2 -- Maximum Path Sum}

The method returns the sum of the maximum simple path connecting two leaves, where a leaf is a node with
no children. If the tree has fewer than two leaves no such path exists and the method returns
\texttt{None}.

\textbf{Observation:} instead of using 0 for this case I return \texttt{None}, because 0 can also be the sum
of a real path (for example two leaves and all keys equal to 0, test
\texttt{zero\_keys\_are\_distinguished\_from\_no\_path}), so the two cases would not be distinguishable.

The method keeps a global maximum, of type \texttt{Option<u32>} and initially \texttt{None}, passed by
mutable reference to a recursive depth-first visit (\texttt{dfs}, line 128) that starts from the root (id 0).
If the node does not exist, \texttt{dfs} returns 0, since we are computing sums. Otherwise it computes
recursively the values of the left and right subtrees (\texttt{l\_max}, \texttt{r\_max}) and the local sum
$\mathit{key} + \mathit{l\_max} + \mathit{r\_max}$.

The local sum is a candidate only if the current node has both children (line 141): a path that bends at
this node goes down to a leaf in the left subtree and to a leaf in the right subtree. With only one child
it would end at the current node, which is not a leaf: see the test
\texttt{path\_cannot\_end\_at\_a\_node\_with\_one\_child}, where the correct answer is 3 and updating at the
root, which has one child, would give 1002. When the node has both children I update the global maximum:
if it is still \texttt{None} it becomes the local sum, otherwise the maximum between the two.

In any case \texttt{dfs} returns $\mathit{key} + \max(\mathit{l\_max}, \mathit{r\_max})$. This is the
invariant of the visit: the returned value is the maximum sum of a path that starts at the node and goes
down to a leaf of its subtree. For a leaf, both subtrees are empty and the value is just its key. When a
node has only one child, the missing side is worth 0 and the other one is at least 0, since keys are
\texttt{u32}, so the maximum always continues towards the existing child.

\textbf{Complexity:}
\begin{itemize}
\item $O(n)$ time, since each node is visited once with constant work;
\item $O(h)$ extra space for the recursion stack, in addition to the tree itself.
\end{itemize}

\section*{Tests}

All 19 tests pass, the code is formatted with \texttt{cargo fmt} and \texttt{cargo clippy} gives no warnings.
\begin{itemize}
\item \texttt{is\_bst}: single node, a valid BST (the tree of the example in the statement), a wrong direct
child on both sides, a grandchild that violates an ancestor, duplicates on the right (accepted) and on the
left (rejected), extreme keys 0 and \texttt{u32::MAX}, key 0 with a left child, a chain of 10\,000 nodes.
\item \texttt{sum\_max\_simple\_path}: single node and chain (\texttt{None}), root with two leaves, a node
with only one child (not an endpoint, and traversed down to the leaf), best path not through the root,
all keys equal to 0 (\texttt{Some(0)}), the tree of the statement.
\item Randomized tests: both methods are compared with brute-force versions on 20\,000 random trees each
(an $O(n^2)$ check of the definition for \texttt{is\_bst}, and for the path all the pairs of leaves through
their lowest common ancestor).
\end{itemize}

\section*{Limitations}
\begin{itemize}
\item Sums are \texttt{u32}, as in the provided \texttt{sum}: a path whose sum is larger than
\texttt{u32::MAX} overflows (panic in debug builds).
\item Both methods are recursive, so on very deep trees the recursion could exceed the stack size.
\end{itemize}

\section*{Use of AI}
I used Claude Code to generate the tests, to review my code, and to translate into English and format this report.

\end{document}
